%==============================================================
% ΚΕΦΑΛΑΙΟ 3 — ΔΕΔΟΜΕΝΑ, ΜΟΝΤΕΛΑ & ΜΕΘΟΔΟΛΟΓΙΑ
%==============================================================
\chapter{Δεδομένα, Μοντέλα και Μεθοδολογία Επικύρωσης}
\label{ch:methodology}

Στο κεφάλαιο 3 θα γίνει αναλυτική περιγραφή των δεδομένων, των μοντέλων 
τεχνητής νοημοσύνης και τελικά της μεθοδολογίας ελέγχου τους.
Στην παρούσα εργασία έγινε χρήση της βάσης FICO HELOC η οποία περιγράφεται 
αναλυτικά στην ενότητα~\ref{sec:dataset}. Η επιλογή της βάσης έγινε κυρίως 
λόγω της κυρίαρχης θέσης της στην βιβλιογραφία όσων αφορά την ερμηνεία και 
παρακολούθηση στον πιστωτικό κίνδυνο. Στην βάση γίνεται προεπεξεργασία που 
αποτελείται από καθαρισμό δεδομένων αλλά και χωρισμό (split) και προσαρμογή 
μεγέθους (scaling) διαδικασίες που περιγράφονται αναλυτικά στην 
ενότητα~\ref{sec:preprocessing}. Η επιλογή των μοντέλων εκτίμησης του 
πιστωτικού κινδύνου ακολουθεί την ίδια λογική. Η Logistic Regression αποτελεί 
σταθερά στον χώρο ενώ οι XGBoost, FFNN τις κύριες νέες εναλλακτικές οι οποίες 
όμως διαφέρουν ως προς την πολυπλοκότητα. Η διαφορά αυτή στο επίπεδο 
πολυπλοκότητας δίνει ένα επιπλέον ενδιαφέρον στην σύγκριση όχι μόνο των 
αποτελεσμάτων τους αλλά κυρίως στις ανάγκες παρακολούθησης και ερμηνείας των 
αποφάσεων τους. Οι περισσότεροι έλεγχοι είναι κοινοί για όλα τα μοντέλα 
ωστόσο υπάρχουν και κάποιοι που γίνονται αποκλειστικά στο XGBoost ως βασικό 
μοντέλο της εργασίας.

\section{Περιγραφή Dataset --- FICO HELOC}
\label{sec:dataset}

Η βάση δεδομένων FICO HELOC δημιουργήθηκε το 2018 από την αμερικανική
εταιρεία FICO (Fair Isaac Corporation) που ειδικεύεται στο data analytics
και συγκεκριμένα στο credit scoring. Η αφορμή της δημιουργίας της ήταν
ένας διαγωνισμός με στόχο την δημιουργία ερμηνεύσιμων μοντέλων
πιστωτικής βαθμολόγησης. Το HELOC (Home Equity Line of Credit) είναι
ένας συγκεκριμένος τύπος δανείου το οποίο έχει ως κάλυψη για την
τράπεζα μέρος της ακίνητης περιουσίας του δανειολήπτη.

Η βάση δεδομένων περιέχει στοιχεία από 10.459 αιτήσεις δανείου. Για
κάθε μια εγγραφή υπάρχουν 23 αριθμητικά χαρακτηριστικά και η βασική
μεταβλητή κριτήριο \textit{RiskPerformance} που παίρνει τιμή
\textit{Good} ή \textit{Bad}, ενώ δεν υπάρχουν προσωπικά στοιχεία όπως
φύλο, ηλικία ή εθνικότητα. Η τιμή της μεταβλητής \textit{RiskPerformance}
εξαρτάται από την αθέτηση ή μη της πληρωμής του δανείου μέσα σε ένα
διάστημα δύο χρόνων από την έκδοσή του. Η αναλογία των αιτήσεων είναι
ομαλή με 52\% αυτών να χαρακτηρίζονται ως \textit{Bad} και 48\% ως
\textit{Good}, όπως φαίνεται στο Σχήμα~\ref{fig:target_dist}.

% -------------------------------------------------------------------
% Σχήμα 3.1 — Κατανομή target μεταβλητής (NB01)
% Τοποθέτηση: αμέσως μετά την αναφορά στην αναλογία 52%/48%
% -------------------------------------------------------------------
\begin{figure}[htbp]
    \centering
    \includegraphics[width=0.55\textwidth]{01_target_distribution}
    \caption[Κατανομή μεταβλητής RiskPerformance]{Κατανομή της μεταβλητής \textit{RiskPerformance} μετά τον
             καθαρισμό της βάσης (9.861 εγγραφές). Η σχεδόν ισορροπημένη
             αναλογία Bad/Good δεν απαιτεί τεχνικές εξισορρόπησης κλάσεων.}
    \label{fig:target_dist}
\end{figure}

Τα αριθμητικά χαρακτηριστικά περιγράφουν το πιστωτικό ιστορικό του
κάθε δανειολήπτη. Ενδεικτικά, \textit{MSinceOldestTradeOpen} είναι οι
μήνες από την παλαιότερη πιστωτική συναλλαγή και
\textit{NetFractionRevolvingBurden} το ποσοστό χρήσης της διαθέσιμης
ανακυκλούμενης πίστωσης, όπως πιστωτικές κάρτες. Για παράδειγμα, σε
μια πιστωτική κάρτα με όριο τα 1.000 ευρώ, εάν έχει γίνει χρήση των
800 ευρώ, τότε το ποσοστό θα είναι 80\%. Όλα τα χαρακτηριστικά είναι
προϊόν επεξεργασίας της FICO με αποτέλεσμα οι σχέσεις μεταξύ τους να
είναι περισσότερο γραμμικές από το σύνηθες, γεγονός που τεκμηριώνεται
οπτικά στον πίνακα συσχέτισης του Σχήματος~\ref{fig:corr_matrix}. Το
γεγονός αυτό εξηγεί και τη μικρή διαφορά στην απόδοση μεταξύ της
Logistic Regression και των XGBoost, FFNN, όπως θα φανεί στο
Κεφάλαιο~\ref{ch:results}.

Ιδιαίτερο χαρακτηριστικό της βάσης αποτελεί η ύπαρξη τιμών $-9$, οι
οποίες συμβολίζουν την απουσία πληροφορίας. Συγκεκριμένα, 598 εγγραφές
αποτελούνται σχεδόν αποκλειστικά από τιμές $-9$, δηλαδή 598
δανειολήπτες δεν είχαν πιστωτικό ιστορικό. Παρά το ελλιπές ιστορικό,
στις εγγραφές αυτές υπάρχει \textit{RiskPerformance} (331 Bad, 267 Good),
αν και για λόγους που θα εξηγηθούν στην συνέχεια δεν θα
χρησιμοποιηθούν.

Η FICO HELOC είναι ευρέως εδραιωμένη στη βιβλιογραφία για την
ερμηνευσιμότητα της μηχανικής μάθησης, έτσι αποτελεί ιδανική επιλογή
για την παρούσα εργασία \citep{fico2018}. Τα δεδομένα της προέρχονται
από πραγματικές αιτήσεις και αποτελέσματα δανείων του τραπεζικού
συστήματος, γεγονός που προσδίδει πρακτικό χαρακτήρα στην παρούσα
εργασία.

\section{Προεπεξεργασία Δεδομένων}
\label{sec:preprocessing}
Όπως αναφέρθηκε και στην παράγραφο~\ref{sec:dataset} 598 εγγραφές από τις 10.459 
αποτελούνται σχεδόν αποκλειστικά από εγγραφές -9. Στην βάση μας αυτό σημαίνει 
απουσία πιστωτικού ιστορικού για τον υποψήφιο δανειολήπτη. Οι εγγραφές αυτές 
αφαιρέθηκαν από την βάση για δύο βασικούς λόγους. Πρώτον, δεν προσφέρουν κάποια 
πληροφορία από την οποία μπορεί να εκπαιδευτούν τα μοντέλα και δεύτερον το ποσοστό 
εμφάνισης τους είναι αρκετά μικρό ώστε η βάση να διατηρήσει την κλίμακα μεγέθους 
της και χωρίς αυτά. Η επιλογή listwise deletion έναντι imputation τεκμηριώνεται από τη βιβλιογραφία 
του missing data analysis, όπου σε περιπτώσεις που η έλλειψη πληροφορίας είναι 
συστηματική (Missing Not At Random) η εισαγωγή εκτιμώμενων τιμών εισάγει 
συστηματικό σφάλμα στις προβλέψεις \citep{little2019statistical, vanbuuren2018flexible}.
Ως τελικό αποτέλεσμα η βάση δεν περιέχει καμία εγγραφή -9. 
Η μόνη μη αριθμητική μεταβλητή RiskPerformance που έπαιρνε τιμές (\textit{Good}, 
\textit{Bad}) μετατράπηκε σε δυαδική χωρίς να χάνεται πληροφορία. Η αλλαγή αυτή 
είναι αναγκαία για την ορθή λειτουργία των μοντέλων ML. Το 0 πλέον αντιστοιχεί 
σε μη αθέτηση ενώ το 1 σε αθέτηση. Για την εκπαίδευση και αξιολόγηση των μοντέλων 
είναι γνωστό ότι χρειάζονται τρία διαφορετικά σύνολα δεδομένων. Αυτά τα σύνολα 
δημιουργήθηκαν από την επεξεργασμένη βάση FICO HELOC με την χρήση stratification 
(διαστρωμάτωση) ώστε η αναλογία 0/1 (\textit{Good}/\textit{Bad}) να είναι ίδια \citep{kohavi1995study}. 
Τα σύνολα training set (70\%) και validation set (15\%) χρησιμοποιούνται κατά την 
εκπαίδευση με το δεύτερο να έχει κύριο ρόλο στον έλεγχο για overfitting 
(υπερπροσαρμογή) και early stopping (πρόωρη διακοπή εκπαίδευσης). Το τελευταίο 
μέρος της βάσης test set (15\%) δεν εμφανίζεται κατά την εκπαίδευση αλλά γίνεται 
χρήση του κατά την τελική αξιολόγηση.

Για την ορθή χρήση των δεδομένων από τα ML μοντέλα είναι αναγκαίο αυτά να είναι 
όλα στην ίδια κλίμακα. Για να επιτευχθεί αυτό χωρίς την διαρροή πληροφορίας 
(data leakage)\citep{kaufman2012leakage, kapoor2023leakage} από το ένα set στο άλλο έγινε χρήση της τεχνικής StandardScaler 
με fit στο training set και transform σε validation και test sets. Αναλυτικά, fit 
είναι η διαδικασία μέσω της οποίας o αλγόριθμος μαθαίνει από τα δεδομένα και 
υπολογίζει μέση τιμή (mean) και τυπική απόκλιση (std) για κάθε δεδομένο, μια 
δυναμική δηλαδή διαδικασία, ενώ transform είναι η εφαρμογή των προηγούμενων 
ανεξάρτητα με τις τιμές των μεταβλητών. Πρακτικά αυτό διασφαλίζει ότι τα δεδομένα 
παραμένουν απομονωμένα και δεν είναι δυνατή η πρόβλεψη των υπόλοιπων συνόλων 
μέσω του training set. Η διαδικασία που περιγράφηκε γίνεται μια φορά στο πρώτο 
notebook του οποίου τα αποτελέσματα αποθηκεύονται σε αρχεία .npy, .pkl για να 
χρησιμοποιηθούν στο υπόλοιπο της εργασίας
(Σχήματα~\ref{fig:feature_dist_1} ,~\ref{fig:feature_dist_2} και~\ref{fig:corr_matrix}).

% -------------------------------------------------------------------
% Σχήμα 3.2 — Κατανομές 23 χαρακτηριστικών (NB01)
% Τοποθέτηση: τέλος ενότητας 3.2, μετά την αναφορά στα .npy/.pkl
% -------------------------------------------------------------------
\begin{figure}[!htbp]
    \centering
    \includegraphics[width=\textwidth]{02_feature_distributions_part1}
    \caption[Κατανομές χαρακτηριστικών FICO HELOC (1/2)]{Κατανομές των πρώτων 11 χαρακτηριστικών της βάσης FICO HELOC μετά
             τον καθαρισμό. Παρατηρείται ποικιλία μορφών κατανομής
             (κανονική, ασύμμετρη, διακριτή) ανάλογα με τη φύση κάθε
             χαρακτηριστικού (1/2).}
    \label{fig:feature_dist_1}
\end{figure}

\begin{figure}[!htbp]
    \centering
    \includegraphics[width=\textwidth]{02_feature_distributions_part2}
    \caption[Κατανομές χαρακτηριστικών FICO HELOC (2/2)]{Κατανομές των υπόλοιπων 12 χαρακτηριστικών της βάσης FICO HELOC μετά
             τον καθαρισμό (2/2).}
    \label{fig:feature_dist_2}
\end{figure}

% -------------------------------------------------------------------
% Σχήμα 3.3 — Πίνακας συσχέτισης (NB01)
% Τοποθέτηση: αμέσως μετά το Σχήμα 3.2 — τα δύο σχήματα
% προεπεξεργασίας παραμένουν μαζί στο τέλος της ενότητας
% -------------------------------------------------------------------
\begin{figure}[!htbp]
    \centering
    \includegraphics[width=\textwidth]{03_correlation_matrix}
    \caption[Πίνακας συσχέτισης χαρακτηριστικών FICO HELOC]{Πίνακας συσχέτισης (lower triangle) των 23 χαρακτηριστικών
             της βάσης FICO HELOC. Η αυξημένη γραμμικότητα μεταξύ των
             χαρακτηριστικών αποτελεί συνέπεια της επεξεργασίας δεδομένων
             από την FICO και εξηγεί εν μέρει τη μικρή διαφορά απόδοσης
             μεταξύ των μοντέλων.}
    \label{fig:corr_matrix}
\end{figure}

\section{Μοντέλα}
\label{sec:models}

\subsection{Logistic Regression}
\label{sec:lr}
Η πιο διαδεδομένη υλοποίηση της Logistic Regression είναι αυτή της βιβλιοθήκης 
scikit-learn σε γλώσσα Python οπότε επιλέχθηκε για την εργασία. Έγινε εφαρμογή 
της με όλες τις default (προεπιλεγμένες) επιλογές παραμέτρων αλλά η σύγκλιση 
δεν ήταν ικανοποιητική. Για τον λόγο αυτό εκπαιδεύτηκε με \texttt{max\_iter=1000} 
αντί 100. Σε κάθε ερευνητική εργασία είναι απαραίτητη η ύπαρξη ενός στόχου, 
ορίου το οποίο να κάνει την αξιολόγηση των αποτελεσμάτων προφανή. Αυτός είναι 
ο ρόλος της LR. Αποτελεί το κύριο μοντέλο στον τραπεζικό τομέα με την πλήρη 
αποδοχή των εποπτικών αρχών λόγω της πλήρους ερμηνευσιμότητάς της, άρα η χρήση 
κάποιου άλλου μοντέλου δικαιολογείται μόνο εάν αυτό χαρακτηρίζεται από καλύτερη 
απόδοση. Τα υπόλοιπα μοντέλα λοιπόν πρέπει να αποδείξουν ότι μπορούν να κάνουν 
καλύτερες προβλέψεις, να περάσουν δηλαδή το όριο που θέτει η LR για να 
δικαιολογηθούν οι επιπλέον δυσκολίες απόρροια της πολυπλοκότητάς τους.


\subsection{XGBoost}
\label{sec:xgboost}

Το μοντέλο XGBoost με το πλήρες όνομα eXtreme Gradient Boosting ανήκει στην
οικογένεια αλγορίθμων τεχνητής νοημοσύνης που χρησιμοποιούν δέντρα αποφάσεων.
Συγκεκριμένα, είναι υλοποίηση του αλγορίθμου Gradient Boosting. Στην πράξη, το
μοντέλο κατά την εκπαίδευσή του φτιάχνει πολλαπλά δέντρα αποφάσεων μικρού βάθους,
τα οποία χρησιμοποιεί στη συνέχεια για να κάνει προβλέψεις. Η δενδρική του αυτή
δομή είναι και ο λόγος που έγινε η επιλογή χρήσης μη κανονικοποιημένων
χαρακτηριστικών (raw features) αντί κλιμακωμένων (scaled), 
καθώς τα δέντρα αποφάσεων είναι αναλλοίωτα ως προς μονότονους μετασχηματισμούς 
των χαρακτηριστικών \citep{hastie2009}. Οι παράμετροι
εκπαίδευσης παρουσιάζονται στον Πίνακα~\ref{tab:xgb_params} ομαδοποιημένοι σε
τρεις κατηγορίες. Αποφεύχθηκε η βελτιστοποίηση υπερπαραμέτρων (hyperparameter
tuning) και οι παράμετροι επιλέχθηκαν ως λογικές αρχικές τιμές που αποφεύγουν
το overfitting \citep{probst2019tunability, bentejac2021comparative}. 
Οι βασικοί λόγοι επιλογής του XGBoost είναι η κυρίαρχη θέση του
στη βιβλιογραφία και η ικανότητα ερμηνευσιμότητάς του μέσω του TreeExplainer του
SHAP \citep{chen2016xgboost}. Όσον αφορά τα μοντέλα μηχανικής μάθησης στον
πιστωτικό κίνδυνο, η υλοποίηση XGBoost χαίρει ευρείας αποδοχής τόσο στην
αποδοτικότητα όσο και στην επικύρωση και παρακολούθησή του. Την ίδια αποδοχή και
αξιοπιστία χαίρει και το TreeExplainer, το οποίο χρησιμοποιείται για να εξηγήσει
τη συμπεριφορά παρόμοιων μοντέλων, και εφόσον είναι συμβατό με το XGBoost,
αποτελεί έναν επιπλέον λόγο επιλογής του. Ως κύριο μοντέλο της εργασίας, πάνω
του θα πραγματοποιηθούν πλήρης ανάλυση SHAP, ανίχνευση μεροληψίας (bias
detection) και παρακολούθηση (monitoring), ενώ η σύγκριση με τη Logistic
Regression θα αποκαλύψει τι επιπλέον απαιτήσεις επιφέρει η αυξημένη πολυπλοκότητα
και πώς αυτές μπορούν να ικανοποιηθούν.

\begin{table}[htbp]
\centering
\caption{Παράμετροι εκπαίδευσης XGBoost}
\label{tab:xgb_params}
\begin{tabular}{llp{7cm}}
\toprule
\textbf{Κατηγορία} & \textbf{Παράμετρος} & \textbf{Τιμή \& Αιτιολόγηση} \\
\midrule
\multirow{3}{*}{Δομή \& Εκπαίδευση}
  & \texttt{n\_estimators} & 300 — αρκετά δέντρα για ακρίβεια χωρίς υπερβολικό υπολογιστικό κόστος \\
  & \texttt{max\_depth}    & 4 — περιορισμός βάθους για αποφυγή overfitting \\
  & \texttt{learning\_rate}& 0.05 — μικρό βήμα μάθησης, συνδυάζεται με τα 300 δέντρα \\
\midrule
\multirow{2}{*}{Τυχαιοποίηση}
  & \texttt{subsample}        & 0.8 — κάθε δέντρο εκπαιδεύεται σε τυχαίο 80\% των εγγραφών \\
  & \texttt{colsample\_bytree}& 0.8 — κάθε δέντρο χρησιμοποιεί τυχαίο 80\% των χαρακτηριστικών \\
\midrule
\multirow{3}{*}{Τεχνικές}
  & \texttt{eval\_metric}        & auc — κατάλληλο για δυαδική ταξινόμηση \\
  & \texttt{random\_state}       & 42 — αναπαραγωγιμότητα αποτελεσμάτων \\
  & \texttt{early\_stopping\_rounds}& 20 - διακοπή εκπαίδευσης εάν η απόδοση στο validation set δεν βελτιωθεί για 20 συνεχόμενους γύρους, αποτρέποντας την υπερπροσαρμογή\\ 
\bottomrule
\end{tabular}
\end{table}

\subsection{Feed-Forward Neural Network}
\label{sec:ffnn}

Μια διαφορετική κατηγορία μοντέλων μηχανικής μάθησης είναι τα νευρωνικά δίκτυα \citep{goodfellow2016deep}.
Στην κατηγορία αυτή ανήκει το μοντέλο FFNN (Feed Forward Neural Network) με την
ιδιαιτερότητα ότι η πληροφορία σε αυτό κινείται μόνο προς μια κατεύθυνση, από
την είσοδο στα κρυφά επίπεδα (hidden layers) και τέλος στην έξοδο. Αναλυτικά
στην εργασία το μοντέλο έχει δύο hidden layers όπου σε κάθε ένα μειώνονται οι
νευρώνες, από 64 σε 32. Η επιλογή αυτή έγινε λόγω του πλήθους των features,
23 σε αριθμό, και έχοντας ως στόχο την αποφυγή της υπερσυμπίεσης της πληροφορίας
μέσω απότομης μείωσης των νευρώνων. Κάθε νευρώνας έχει και ένα βάρος με το οποίο
φιλτράρει την πληροφορία σε μη δυαδική μορφή και με όρια διαφορετικά του $(0,1)$,
για τον λόγο αυτό στην έξοδο εφαρμόζεται μια sigmoid activation συνάρτηση που
μετατρέπει την πληροφορία σε πιθανότητα ανάμεσα σε 0 και 1, μορφή κατάλληλη για
την εργασία μετά την εφαρμογή ορίου μετατροπής (threshold). Κάθε νευρώνας έχει
δύο καταστάσεις «ενεργοποιημένος» και «μη ενεργοποιημένος». Η συνάρτηση που
αποφασίζει για την κατάσταση αυτή είναι η ReLU activation \citep{glorot2011deep} η οποία επιστρέφει
μηδέν σε αρνητικές τιμές «μη ενεργοποίηση» και τις τιμές ατόφιες για θετικές
για κατάσταση «ενεργοποίησης». Για την εκπαίδευσή του το μοντέλο απαιτεί
κανονικοποιημένα δεδομένα, οπότε σε αντίθεση με το XGBoost έγινε επεξεργασία
των δεδομένων με προσοχή στο data leakage όπως αναφέρθηκε στην
ενότητα~\ref{sec:preprocessing}. Επιπλέον, για την αποφυγή overfitting στα
δεδομένα ορίστηκε Dropout=0.3/0.2 \citep{srivastava2014dropout} το οποίο απενεργοποιεί τυχαία 30\% των νευρώνων
έτσι ώστε το μοντέλο να βασίζεται σε όλους τους νευρώνες του σχεδόν ισότιμα.
Η μεταβλητή Epochs=100 με EarlyStopping(patience=15) \citep{prechelt1998early} ορίζει το πόσες φορές θα
εκπαιδευτεί το μοντέλο στα δεδομένα με όριο τις 100 φορές, εκτός εάν σε 10
συνεχόμενες εποχές δεν βελτιωθεί το αποτέλεσμά του. Η μεταβλητή αυτή είναι
σημαντική γιατί επιτρέπει να μην οριστεί ακριβές όριο μετεκπαίδευσης που θα
μπορούσε να οδηγήσει σε χειρότερο αποτέλεσμα αλλά και υπερεκπαίδευση. Οι
υπόλοιπες τεχνικές παράμετροι εκπαίδευσης παρουσιάζονται συνοπτικά στον
Πίνακα~\ref{tab:ffnn_params}. Τα δεδομένα της HELOC λόγω της επεξεργασίας που
έχει γίνει από την FICO έχουν αυξημένη γραμμικότητα μεταξύ τους. Για τον λόγο
αυτό οι διαφορές ανάμεσα στην απόδοση των μοντέλων θα είναι μειωμένες, γεγονός
που δεν αποτελεί μειονέκτημα για την εργασία καθώς ο σκοπός είναι η σύγκριση
των μοντέλων ως προς τις ανάγκες παρακολούθησης και ερμηνείας τους. Η σύγκριση
της FFNN με τα υπόλοιπα μοντέλα θα γίνει κυρίως πάνω σε αυτό τον σκοπό.

\begin{table}[htbp]
\centering
\caption{Παράμετροι εκπαίδευσης Feed-Forward Neural Network}
\label{tab:ffnn_params}
\begin{tabular}{llp{7cm}}
\toprule
\textbf{Κατηγορία} & \textbf{Παράμετρος} & \textbf{Τιμή \& Αιτιολόγηση} \\
\midrule
\multirow{2}{*}{Βελτιστοποίηση}
  & \texttt{optimizer}      & Adam \citep{kingma2015adam} --- προσαρμοστικός αλγόριθμος, standard επιλογή για FFNN \\
  & \texttt{learning\_rate} & 0.001 --- προεπιλογή Adam, σταθερή σύγκλιση \\
\midrule
Συνάρτηση κόστους
  & \texttt{loss}           & binary crossentropy --- θεωρητικά ορθή για δυαδική ταξινόμηση \citep{goodfellow2016deep} \\
\midrule
Εκπαίδευση
  & \texttt{batch\_size}    & 64  --- ισορροπία ταχύτητας και σταθερότητας εκπαίδευσης \\
\bottomrule
\end{tabular}
\end{table}

\section{Μεθοδολογία Επικύρωσης}
\label{sec:val_methodology}

\subsection{Αξιολόγηση Απόδοσης}
\label{sec:performance}
Τα μοντέλα αφού εκπαιδευτήκαν στο train και validation δοκιμάστηκαν στο test set.
Για την αξιολόγηση των μοντέλων χρησιμοποιήθηκε πλήθος δεικτών που αποτελούν 
standard στον τομέα του credit scoring \citep{hand1997statistical, anderson2007credit}.


\paragraph{AUC-ROC \& Gini.}
Ο δείκτης AUC-ROC μετράει το εμβαδόν κάτω από την καμπύλη ROC και αποτελεί τον
σημαντικότερο δείκτη όσων αφορά τον πιστωτικό κίνδυνο \citep{hand2009measuring}. Η καμπύλη ROC σχηματίζεται
στον άξονα True Positive Rate (ποσοστό Bad που εντοπίστηκαν σωστά) και False
Positive Rate (ποσοστό Good που κατατάχθηκαν λανθασμένα ως Bad) και το εμβαδόν
μετράει την ικανότητα του μοντέλου να κάνει σωστές προβλέψεις σε κάθε threshold,
δηλαδή σε κάθε όριο πιθανότητας πάνω από το οποίο η πρόβλεψη μετατρέπεται σε Bad.
Η AUC-ROC παίρνει τιμές κοντά στο 0,5 για μοντέλα χωρίς προβλεπτική ικανότητα
και έτσι τα αποτελέσματα των ικανών για προβλέψεις μοντέλων συνήθως κυμαίνονται
στο πεδίο $(0{,}5,\ 1)$. Για πιο ξεκάθαρα αποτελέσματα στην κλίμακα $(0,1)$
χρησιμοποιείται ο δείκτης Gini ο οποίος προκύπτει από την σχέση $2 \times
\text{AUC} - 1$.

\paragraph{KS (Kolmogorov-Smirnov).}
Ενώ οι προηγούμενοι δείκτες φανερώνουν την απόδοση ανεξάρτητα από threshold,
ο KS εντοπίζει σε ποια τιμή του το μοντέλο έχει την βέλτιστη απόδοση. Για να
το πετύχει αυτό δημιουργεί μια φθίνουσα λίστα επικινδυνότητας και δύο αθροιστικές
κατανομές. Οι κατανομές οπτικοποιούνται ως καμπύλες και φανερώνουν το ποσοστό
Bad και Good μέχρι κάθε συγκεκριμένο σημείο της λίστας. Το σημείο στο οποίο οι
καμπύλες διαφέρουν το μέγιστο είναι το κατάλληλο threshold.

\paragraph{Πίνακας Σύγχυσης \& Ανάλυση Ορίου Απόφασης.}
Επειδή το κόστος κάθε Ψευδώς Θετικού (False Positive) δεν είναι πάντα ίδιο
με το κόστος ενός Ψευδώς Αρνητικού (False Negative) χρειάζονται και άλλες μετρικές
όπως ο Πίνακας Σύγχυσης (Confusion Matrix). Για συγκεκριμένο threshold (0.5 στην
παρούσα εργασία) ο Πίνακας Σύγχυσης είναι ένας $2\times2$ πίνακας με το σύνολο
των TP (True Positive), TN (True Negative), FP (False Positive) και FN (False
Negative). Στον πιστωτικό κίνδυνο η ελαχιστοποίηση του FN, των δανειοληπτών που
δεν θα μπορέσουν να πληρώσουν το δάνειο, είναι απαραίτητη. Συμπληρωματικά
λειτουργεί η Ανάλυση Ορίου Απόφασης (Threshold Analysis). Το Threshold Analysis
αποτελείται από γραφήματα τριών μετρικών, της Precision, Recall και F1 οι οποίες
υπολογίζονται για κάθε threshold. Η Precision υπολογίζει από τους προβλεπόμενους
Bad πόσοι ήταν πραγματικά Bad, η Recall υπολογίζει από τους πραγματικούς Bad
πόσους εντόπισε το μοντέλο, ενώ η F1 είναι ο αρμονικός μέσος των προηγούμενων
δύο.

\paragraph{Ανάλυση Κατά Δεκατημόρια (Decile Analysis).}
Κλασική μετρική στον οικονομικό τομέα, χωρίζει το δείγμα ανάλογα με την
προβλεπόμενη επικινδυνότητα σε 10 ομάδες και ελέγχει το ποσοστό προβλέψεων Bad
(Bad Rate) σε αυτές. Ένα καλό προβλεπτικό μοντέλο αναμένεται να έχει μεγάλο
Bad Rate στα πρώτα Decile ενώ μικρότερο στα επόμενα. Στην εργασία χρησιμοποιείται
επίσης το Αθροιστικό Ποσοστό Αθέτησης (Cumulative Bad Rate) στο οποίο φαίνεται
το αθροιστικό ποσοστό των Bad ανά δεκατημόριο και συγκρίνεται με την διαγώνιο
που αντιστοιχεί σε μοντέλο χωρίς προβλεπτική ικανότητα.

\paragraph{Βαθμονόμηση (Calibration).}
Η Βαθμονόμηση εξετάζει κατά πόσο οι προβλεπόμενες πιθανότητες αθέτησης
αντιστοιχούν στα πραγματικά δεδομένα. Η μέτρηση αυτή είναι αναγκαία στον
πιστωτικό κίνδυνο. Οι κανόνες της Basel III επιβάλλουν την ακριβή τιμολόγηση
επιτοκίων και τον υπολογισμό αποθεματικών κεφαλαίου τα οποία βασίζονται στις
προβλέψεις των ποσοστών αθέτησης \citep{basel2017}. Για να κρίνουμε πόσο ακριβές είναι αυτό το
ποσοστό, ομαδοποιούμε τους δανειολήπτες με βάση αυτή την πιθανότητα και για
κάθε ομάδα (bin) της μορφής $(x,\ x+\Delta y)\%$ βρίσκουμε την πραγματική
πιθανότητα αθέτησης. Έτσι στο Calibration Plot ένα μοντέλο με τέλεια ικανότητα
βαθμονόμησης θα ακολουθεί την διαγώνιο εφόσον η ομαδοποίηση θα ταυτίζεται με
τα πραγματικά αποτελέσματα. Ο στόχος των μοντέλων είναι να πλησιάζουν όσο
περισσότερο την διαγώνιο. Η μετρική αυτή διαφέρει από την AUC καθώς η δεύτερη
αρκείται στην ταξινόμηση των δανειοληπτών και όχι στην ακριβή πρόβλεψη του
ποσοστού αθέτησης.

\subsection{Ανίχνευση Μεροληψίας}
\label{sec:bias_methodology}

Η Ανίχνευση Μεροληψίας (Bias Detection) ελέγχει εάν οι προβλέψεις του μοντέλου
χειρίζονται κάθε υποομάδα με τον ίδιο τρόπο. Ο διαχωρισμός των ομάδων γίνεται
συνήθως με βάση το φύλο, ηλικία ή εθνικότητα του δανειολήπτη. Στην παρούσα
εργασία επειδή η βάση δεν περιέχει προσωπικά δεδομένα οι ομάδες χωρίζονται με
βάση την μεταβλητή \texttt{AverageMInFile}. Η χρήση proxy μεταβλητής αποτελεί καθιερωμένη πρακτική στη βιβλιογραφία fairness 
όταν τα προστατευόμενα χαρακτηριστικά δεν είναι διαθέσιμα στη βάση δεδομένων 
\citep{chen2019fairness}, και επιτρέπει την εμπειρική επίδειξη της μεθοδολογίας 
ανίχνευσης μεροληψίας χωρίς να αντικαθιστά την ανάλυση σε πραγματικά 
δημογραφικά χαρακτηριστικά. Η ομάδα Α με
\texttt{AverageMInFile} κάτω της διαμέσου έχει μικρότερο πιστωτικό
παρελθόν ενώ η Β με τιμές πάνω της διαμέσου μεγαλύτερο. Η σύγκριση των ποσοστών
έγκρισης κάθε ομάδας είναι η κύρια διαφορά με την AUC και Calibration οι οποίες
κάνουν μόνο ατομικούς ελέγχους σε κάθε δανειολήπτη. Οι δύο κυριότερες μετρικές
της Bias Detection είναι η Δημογραφική Ισότητα (Demographic Parity) και η
Ανισομερής Επίπτωση (Disparate Impact). Στην DPD (Demographic Parity Difference)
γίνεται σύγκριση του ποσοστού έγκρισης των δύο ομάδων με αποδεκτό όριο το
$0{,}1$. Στην DI (Disparate Impact Ratio) υπολογίζεται ο λόγος των ποσοστών
έγκρισης και το όριο είναι το $0{,}8$. Οι δύο αυτές μετρικές είναι
συμπληρωματικές και δίνουν διαφορετικές πληροφορίες. Στην εργασία παρατηρείται
bias, δηλαδή οι τιμές ξεπερνούν τα όρια που αναφέρθηκαν, οπότε εφαρμόζονται δύο τεχνικές μετριασμού που ανήκουν στην καθιερωμένη ταξονομία 
pre/in/post-processing \citep{kamiran2012data}: \texttt{ExponentiatedGradient} 
\citep{agarwal2018reductions} με περιορισμό Demographic Parity για την Logistic 
Regression (in-processing) και χειροκίνητος ThresholdOptimizer \citep{hardt2016equality} 
μέσω grid search για το XGBoost (post-processing).
Και οι δύο τεχνικές προσαρμόζουν τα όρια απόφασης ξεχωριστά για κάθε ομάδα.
Με τον τρόπο αυτό τα μοντέλα έχουν αποδεκτή, μικρότερη μεροληψία αλλά η
απόδοσή τους επηρεάζεται --- για την Logistic Regression μειώνεται το AUC,
ενώ για το XGBoost το AUC παραμένει αμετάβλητο και μειώνεται το Balanced
Accuracy. Το αποτέλεσμα αυτό είναι αναμενόμενο καθώς η διαδικασία στοχεύει
στην διατήρηση μικρής διαφοράς ανάμεσα στα ποσοστά έγκρισης για κάθε ομάδα
και όχι στην καλύτερη δυνατή πρόβλεψη για κάθε δανειολήπτη ξεχωριστά.

\subsection{Σταθερότητα Πληθυσμού (PSI)}
\label{sec:psi}

Η Σταθερότητα Πληθυσμού (PSI) υπολογίζει εάν τα δεδομένα εκπαίδευσης
αποκλίνουν από τα δεδομένα πρόβλεψης της τρέχουσας χρονικής περιόδου.
Η μετρική αυτή στοχεύει στην ανίχνευση αλλαγών στα δεδομένα εισόδου
και όχι στην αξιολόγηση της ποιότητας των προβλέψεων. Στην πράξη
χωρίζει τις προβλεπόμενες πιθανότητες σε δέκα ομάδες με ίσα διαστήματα
από μηδέν μέχρι ένα. Έπειτα για κάθε ομάδα υπολογίζεται το ποσοστό
εγγραφών για τα δύο σύνολα (reference dataset και current dataset).
Τέλος με τον μαθηματικό τύπο:
\begin{equation}
    PSI = \sum_{i=1}^{n} (A_i - E_i) \times \ln\left(\frac{A_i}{E_i}\right)
    \label{eq:psi}
\end{equation}
υπολογίζεται ο δείκτης PSI, όπου $A_i$ είναι το ποσοστό εγγραφών στο
current dataset και $E_i$ στο reference dataset για κάθε ομάδα $i$.
Εάν ο δείκτης είναι μικρότερος του $0{,}10$ έχουμε σταθερή κατανομή,
εάν είναι ανάμεσα σε $0{,}10$ και $0{,}25$ υπάρχει μέτρια μετατόπιση
στα δεδομένα και η παρακολούθηση κρίνεται αναγκαία, ενώ για τιμές
μεγαλύτερες του $0{,}25$ απαιτείται επανεκπαίδευση του μοντέλου.
Στην εργασία τα δεδομένα της βάσης προέρχονται από μια συγκεκριμένη
χρονική περίοδο οπότε έγινε προσομοίωση μελλοντικών δεδομένων με την
χρήση Γκαουσιανού θορύβου. Έγινε η υπόθεση ότι η μελλοντική οικονομία
θα υποστεί μια ύφεση με τους βασικούς δείκτες να αλλάζουν.
Συγκεκριμένα, οι δείκτες \texttt{ExternalRiskEstimate} και
\texttt{AverageMInFile} μειώθηκαν διότι κατά την διάρκεια μιας
οικονομικής ύφεσης το πιστωτικό προφίλ χειροτερεύει και το πιστωτικό
ιστορικό ελαττώνεται, ενώ αυξάνεται το χρέος και κατά συνέπεια ο
δείκτης \texttt{NetFractionRevolvingBurden}. Έγινε προσομοίωση τριών εντάσεων, ήπιας, μέτριας και σοβαρής επιδείνωσης
με κλιμακώμενη ένταση θορύβου ανάλογα με τη σοβαρότητα. Με τον τρόπο αυτό γίνεται
ευδιάκριτη η χρησιμότητα της μετρικής.

\section{Εξειδικευμένη Ανάλυση XGBoost}
\label{sec:xgboost_analysis}

Τα μοντέλα τεχνητής νοημοσύνης συνήθως παράγουν καλύτερα αποτελέσματα στον
πιστωτικό κίνδυνο εισάγοντας όμως παράλληλα και έναν συμβιβασμό (trade-off).
Όσο πιο περίπλοκα είναι τόσο λιγότερο ερμηνεύσιμη είναι η διαδικασία που
ακολούθησαν για να εξάγουν το αποτέλεσμα. Το μοντέλο XGBoost είναι το κύριο
μοντέλο της εργασίας αλλά και το κυρίαρχο μοντέλο στον τομέα του διότι
βρίσκεται στο σημείο της χρυσής τομής του trade-off. Το XGBoost προσφέρει
αισθητή βελτίωση αποδοτικότητας έναντι της Logistic Regression, διατηρώντας
παράλληλα επαρκή ερμηνευσιμότητα για την ανάλυση που ακολουθεί. Σε αυτή την
ενότητα θα γίνει ανάλυση των μετρικών και δεικτών που εξηγούν τους λόγους για
τους οποίους το μοντέλο πήρε τις συγκεκριμένες αποφάσεις αλλά και των τεχνικών
παρακολούθησης των μοντέλων μετά το τέλος της εκπαίδευσής τους. Οι μετρήσεις
αυτές δεν αποτελούν αυτοσκοπό αλλά είναι αναγκαίες καθώς αποτελούν
προαπαιτούμενο από τις διάφορες ρυθμιστικές αρχές (Basel, EBA) \citep{basel2017, eba2023}.

\subsection{Ερμηνευσιμότητα --- SHAP}
\label{sec:shap}

Για να εξηγήσουμε μια απόφαση έγκρισης δανείου δεν αρκεί ένα ποσοστό.
Η χρήση εργαλείων που εξηγούν την πιθανότητα αθέτησης και κατά συνέπεια
την απόφαση έγκρισης ή μη του δανείου είναι αναγκαία. Το SHAP (SHapley
Additive exPlanations) μέσω της θεωρίας Shapley values του τομέα της
θεωρίας παιγνίων καταφέρνει να αποκωδικοποιήσει τη συμπεριφορά του
μοντέλου. Μέσω του SHAP μπορούμε να καταλάβουμε κατά πόσο επηρέασε η
κάθε μεταβλητή το αποτέλεσμα. Η λειτουργία του είναι διπλή, διακρίνεται
στην ολική ανάλυση των δεδομένων (Global SHAP) και στην ατομική ανάλυση
(Local SHAP).Οι τιμές αυτές είναι αθροιστικές: το άθροισμά τους μαζί με
την τιμή αναφοράς (base value) ισούται με την τελική πρόβλεψη του μοντέλου
για τον συγκεκριμένο δανειολήπτη.
Το Local SHAP υπολογίζει για κάθε χαρακτηριστικό του
δανειολήπτη μια τιμή SHAP που συμβολίζει τη συνεισφορά του στην τελική
πρόβλεψη. Από τις τιμές αυτές γίνεται δυνατή η ικανοποίηση του
«δικαιώματος εξήγησης» του αιτούντος πελάτη. Το Global SHAP υπολογίζει
τη μέση απόλυτη τιμή από κάθε SHAP τιμή και παράγει πιο γενικά
αποτελέσματα που εξηγούν τη λειτουργία του μοντέλου. Από τα αποτελέσματα
γίνονται σαφή τα χαρακτηριστικά τα οποία αυτό θεωρεί σημαντικότερα.

Η λειτουργία του PDP (Partial Dependence Plots) είναι συμπληρωματική του SHAP.
Ενώ το SHAP απαντάει στο πόσο επηρεάζει κάθε χαρακτηριστικό, το PDP
συμπληρώνει τον τρόπο με τον οποίο το κάνει. Για κάθε χαρακτηριστικό
δοκιμάζει πολλαπλές τιμές του σε όλους τους υποψήφιους δανειολήπτες και
καταγράφει τις αλλαγές στην τελική πρόβλεψη. Ως αποτέλεσμα, παράγονται
διαγράμματα που παρουσιάζουν τη μη γραμμική συσχέτιση των χαρακτηριστικών
με τον υπολογισμό του ποσοστού αθέτησης. Η μη γραμμική ``σκέψη'' του
μοντέλου είναι αυτή που το ξεχωρίζει από την Logistic Regression.

\subsection{Σταθερότητα Χαρακτηριστικών (CSI)}
\label{sec:csi}

Ο δείκτης CSI (Characteristic Stability Index) λειτουργεί συμπληρωματικά
του PSI. Χρησιμοποιεί τον ίδιο μαθηματικό τύπο (βλ.\ εξίσωση~\ref{eq:psi})
για κάθε χαρακτηριστικό αντί μόνο του μεμονωμένου ποσοστού πρόβλεψης.
Οι τιμές CSI κάτω του $0{,}10$ μεταφράζονται ως σταθερό χαρακτηριστικό,
οι τιμές μεταξύ $0{,}10$ και $0{,}25$ ως ύπαρξη μέτριας μετατόπισης,
ενώ για τιμές άνω του $0{,}25$ ύπαρξη σημαντικής αστάθειας. Σημαντική
πληροφορία αποτελεί ότι ο δείκτης CSI βασίζεται σε ιστόγραμμα (histogram)
οπότε είναι ιδιαίτερα ευαίσθητος όταν αξιολογεί διακριτές τιμές. Ακόμα
και μικρές μεταβολές των διακριτών τιμών μπορούν να προκαλέσουν τεχνητά
υψηλές τιμές του δείκτη. Στην εργασία η χρησιμότητα αυτού του δείκτη
θα φανεί μέσω του τροποποιημένου dataset.


\subsection{Παρακολούθηση Μοντέλου}
\label{sec:monitoring_methodology}

Έπειτα από το πέρας της εκπαίδευσης, της επικύρωσης και των τελικών
ελέγχων η τελευταία ασπίδα προστασίας για τη διαφύλαξη της σωστής
λειτουργίας του μοντέλου είναι η παρακολούθηση μοντέλου (Monitoring).
Η φθορά απόδοσης του μοντέλου με το χρόνο (Model Drift) είναι
αναμενόμενη. Για τον λόγο αυτό ο συνεχής έλεγχος είναι προαπαιτούμενο
και βασική απαίτηση των εποπτικών αρχών \citep{eba2023}. Η παρακολούθηση μοντέλου
εστιάζει στους δείκτες PSI --- παρακολουθεί αλλαγές στα ποσοστά των
προβλέψεων --- CSI --- παρακολουθεί αλλαγές στα ποσοστά των
χαρακτηριστικών --- και AUC over time --- παρακολουθεί την προβλεπτική
ικανότητα του μοντέλου. Στην παρούσα εργασία ο έλεγχος γίνεται σε
δεδομένα που προσομοιώνουν ήπια, μέτρια και σοβαρή οικονομική
επιδείνωση για τα μοντέλα XGBoost και FFNN. Η βασική διαφορά των
μοντέλων κατά την διαδικασία βρίσκεται στην ανάγκη χρήσης της
συνάρτησης KernelExplainer \citep{lundberg2017shap} 
αντί της TreeExplainer \citep{lundberg2020local} για το μοντέλο FFNN,
η οποία αναμένεται να αυξήσει σημαντικά την πολυπλοκότητα και ως
απόρροια τις ώρες υπολογισμού.